\documentclass[11pt]{amsart}
\usepackage[margin=1.1in]{geometry}
\usepackage{amsmath,amssymb,amsthm}
\usepackage{hyperref}
\newtheorem{theorem}{Theorem}
\newtheorem{proposition}[theorem]{Proposition}
\newtheorem{lemma}[theorem]{Lemma}
\newtheorem*{propsharp}{Proposition 2$'$}
\theoremstyle{remark}
\newtheorem{remark}[theorem]{Remark}
\newcommand{\R}{\mathbb{R}}
\newcommand{\eps}{\varepsilon}
\newcommand{\ip}[2]{\langle #1,#2\rangle}

\title{A sharper enclosure of the one-loop shrinker $S_1$}
\author{Supplement to the working notes for W.\,J.~Ogden and M.~Warren}
\date{November 6, 2026. Proof, code and draft by Claude (Anthropic) at M.~Warren's request.}

\begin{document}
\begin{abstract}
The step-1 supplement encloses the conical end of $S_1$ at $\sigma=8$ by its first-order state, with an uncertainty of $3.4\cdot10^{-8}$. That uncertainty was the floor for every later certificate that uses the profile. Here the first-order state is replaced by the exact solution of the linearized end equation, written as a Laplace integral; only a Lyapunov--Perron remainder that is cubic in $\eps$ is left. The uncertainty at $\sigma=8$ drops to $4.4\cdot10^{-13}$. The unchanged CAPD shooting then gives $\eps^*$ to relative accuracy $4\cdot10^{-10}$ and the tip to $10^{-9}$. Downstream, the certified residual of the slow-mode approximant $\hat\psi_3$ drops from $3.8\cdot10^{-3}$ to $1.3\cdot10^{-6}$.
\end{abstract}
\maketitle



Notation is that of the step-1 supplement: $S=8$, $\mu(\sigma)=(\sigma/4)^{1/3}$, $E$, $w$, $g$, the norms $\|\cdot\|_0$, $\|\cdot\|_1$, the constants $c=1/\sqrt3$, $\kappa$, $\nu_0$, $\nu_1$, and the diagonalization $Y=(k,k',k'')^T=P(\sigma)Z$ of $k'''+\frac\sigma4k=N$, with $N=-k(kk'+J)$ and $J=\int_\sigma^\infty k'^2$. As there, $\eps$ is defined by $Z_0(S)=\eps$.

\section{The exact linear end}
\begin{lemma}\label{lem:laplace}
For $\sigma>0$ let $a=(\sigma/4)^{1/3}$ and let $C$ consist of the rays $t=-a+re^{\pm3\pi i/4}$, $r\ge0$, oriented from $\infty e^{-3\pi i/4}$ to $\infty e^{3\pi i/4}$. Put
\[
K_j(\sigma)=\frac1{2\pi i}\int_C t^j e^{\sigma t+t^4}\,dt\qquad(j\ge0).
\]
Then $K_j$ is real, $K_j=K_0^{(j)}$, $K_0'''+\frac\sigma4K_0=0$, and $|K_j(\sigma)|\le\frac1\pi\int_0^\infty(a+r)^je^{-\frac34\sigma a-r^4}dr$. In particular $K_0$ spans the decaying solutions.
\end{lemma}

\begin{proof}
On the rays,
\[
\mathrm{Re}(\sigma t+t^4)=a^4+2\sqrt2a^3r-2\sqrt2ar^3-r^4-\sigma a-\sigma r/\sqrt2\le-\tfrac34\sigma a-r^4,
\]
because $a^3=\sigma/4$ makes the terms linear in $r$ cancel. This gives absolute convergence, the bound, and differentiation under the integral. The contour may be moved locally in $\sigma$ by Cauchy's theorem. Then $K_0'''+\frac\sigma4K_0=\frac1{8\pi i}\int_C\frac d{dt}e^{\sigma t+t^4}dt=0$. The two rays are complex conjugates and the integrand is real on $\R$, so $K_j=\mathrm{Im}(I_j)/\pi$, where $I_j$ is the integral over the upper ray. The bound decays like $e^{-E}$. The other two solutions grow like $e^{E/2}$, so the decaying solutions form a line.
\end{proof}

The values $K_0,K_1,K_2$ at $S$ are computed with Arb's validated quadrature on $r\in[0,8]$ plus the tail bound of the lemma (\texttt{laplace.py}). The check $K_3+\frac S4K_0=O(10^{-49})$ is satisfied. Normalize $\hat Y=c_n(K_0,K_1,K_2)$ so that $\hat Z=P^{-1}\hat Y$ has $\hat Z_0(S)=1$. Then
\[
\hat Y(8)=(0.97939551476\ldots,\ -1.27198455271\ldots,\ 1.60490957236\ldots),
\]
\[
\hat Z_1(8)=-0.0103022+0.00041998\,i .
\]
For comparison, the first-order value $\gamma^{(1)}=-0.0103562+0.000333\,i$ of the step-1 supplement differs from $\hat Z_1(8)$ by $O(\kappa^2)$.

\section{The sharp end state}
\begin{propsharp}
Let $|\eps|\le\eps_{\max}=1.62\cdot10^{-4}$. The conical end of Proposition~2 of the step-1 supplement has, at $\sigma=S=8$,
\[
(k,k',k'')^T=\eps\hat Y(S)+P(S)\,(0,d,\bar d)^T,\qquad |d|\le r_1^{\rm nl}=4.37\cdot10^{-13},
\]
and $p=4k''$, $m=k'$, $q=S+4kk'+4J$, with
\[
J=\eps^2\hat J+J',\qquad \hat J=c_n^2\int_S^\infty K_1^2=0.631426835(1),\qquad |J'|\le3.6\cdot10^{-16}.
\]
\end{propsharp}

\begin{proof}
\emph{Linear part.} $\eps\hat Z$ solves the linear $Z$-system exactly, with $\hat Z_0(S)=1$. By Lemma~\ref{lem:laplace} it is the fixed point of the linear map $T$ of the step-1 proof (the case $N=0$, $\eps=1$). The step-1 inequalities with $N=0$ hold for $r_0^L=3.74\cdot10^{-3}$ and $r_1^L=1.294\cdot10^{-2}$:
\[
r_1^L\ge\kappa(1+r_0^L+r_1^L),\qquad r_0^L\ge\tfrac c2r_1^L .
\]
Hence $\|\hat Z_0-\Phi_0(\cdot,S)\|_0\le r_0^L$ and $\|\hat Z_1\|_1\le r_1^L$. Since $\|\Phi_0(\cdot,S)\|_0=1$, the total weighted size is at most $\hat{\mathcal R}=1+r_0^L+2r_1^L$.

\emph{Cubic remainder.} $Z^{\rm nl}=Z-\eps\hat Z$ is a fixed point of $T$ with $\eps=0$ and with $N$ evaluated at $\eps\hat Y+PZ^{\rm nl}$. Now the coupling acts only on $Z^{\rm nl}$. With $\mathcal R=\eps_{\max}\hat{\mathcal R}+r_0^{\rm nl}+2r_1^{\rm nl}$, the step-1 estimates give the invariance conditions
\[
r_0^{\rm nl}\ge\tfrac c2r_1^{\rm nl}+\nu_0\mathcal R^3,\qquad r_1^{\rm nl}\ge\kappa(r_0^{\rm nl}+r_1^{\rm nl})+\nu_1\mathcal R^3 .
\]
They hold for $r_0^{\rm nl}=8.6\cdot10^{-13}$ and $r_1^{\rm nl}=4.37\cdot10^{-13}$, with contraction constant $L=0.3014$ (\texttt{tail2.py}, Arb at 256 bits). Then $d=Z_1^{\rm nl}(S)$ and $Z^{\rm nl}_0(S)=0$, so $Z_0(S)=\eps$ as in step 1.

\emph{The integral $J$.} $k'=\eps\hat Y_1+k'_{\rm nl}$ with $|k'_{\rm nl}|\le\mu(r_0^{\rm nl}+2r_1^{\rm nl})/w$. As in step 1, $\int_S^\infty\mu^2w^{-2}\le\frac{\mu(S)}2+\frac1{12S}$. So $\|k'_{\rm nl}\|_{L^2(S,\infty)}\le B=(r_0^{\rm nl}+2r_1^{\rm nl})(\frac{\mu(S)}2+\frac1{12S})^{1/2}$, and $|J'|\le2|\eps|\hat J^{1/2}B+B^2$. Since $\mathrm{Re}(t+s)\le-2a<0$ on $C\times C$,
\[
\int_S^\infty K_1^2=-\frac1{(2\pi i)^2}\int_C\!\int_C\frac{ts\,e^{t^4+s^4+S(t+s)}}{t+s}\,dt\,ds .
\]
The double integral is evaluated in Arb on $[0,6]^2$ in the ray parameters (\texttt{jhat.py}). By the estimate in the proof of Lemma~\ref{lem:laplace}, the rest is at most $\frac{4G_0G_R}{4\pi^2a}<5\cdot10^{-573}$, where $G_0=e^{c_0}(a\Gamma(\frac54)+\frac{\sqrt\pi}4)$ and $G_R\le e^{c_0-R^4}(a+R)/(4R^3-\frac1{a+R})$ with $c_0=-\frac34Sa$ (\texttt{jhat\_rig.py}). The result is $\int_8^\infty K_1^2=1.4404409694(1)\cdot10^{-9}$.
\end{proof}

\section{Shooting}
\texttt{gen\_input2.py} builds the outward-rounded parallelepiped of Proposition~2$'$ for an $\eps$-interval. Its directions are $\partial_\eps$ (now carrying $dq/d\eps$ at the centre), $\mathrm{Re}\,d$, $\mathrm{Im}\,d$, $\partial_q$ (the quadratic remainder and $J'$), and $\partial_\theta$. The CAPD programs \texttt{s1proof.cpp} and \texttt{s1curve.cpp} of the step-1 supplement are unchanged.
\begin{itemize}
\item $M(1.45\cdot10^{-4})\in[0.04339455686,\,0.04339455704]$.
\item $M(1.62\cdot10^{-4})\in[-0.03933409755,\,-0.03933409733]$. The step-1 enclosures had width $1.7\cdot10^{-5}$.
\item Bisection on certified signs (\texttt{bracket2.py}) gives
\[\eps^*\in[1.5346198892965,\ 1.5346198899299]\cdot10^{-4}.\]
\item On the bracket set:
  \begin{itemize}
  \item $\tau_{\rm tip}\in[10.249077650575,\,10.249077651520]$;
  \item $\sigma_{\rm tip}=8-\tau_{\rm tip}\in[-2.2490776515195,\,-2.2490776505754]$;
  \item $D\in[2.50080627654,\,2.50080627668]$;
  \item $k(0)\in[1.53982063762,\,1.53982063822]$.
  \end{itemize}
\item The geometric checks (a)--(d) of the step-1 supplement pass unchanged. The double point now lies in $\tilde x\in[-0.83152,-0.82932]$.
\end{itemize}

\section{Consequence for the slow-mode residual}
\texttt{s1prof.cpp} produces 51{,}246 CAPD boxes in $\tau\in[0,\tau_{\rm tip}]$, each with an enclosure at its centre (width $\le2.7\cdot10^{-9}$) and a hull. Write $R=(L-\tilde\lambda)\hat\psi_3=-u''''-a_2u''-a_1u'+a_0u$, where $a_2=k^2$, $a_1=5km+q/4$, $a_0=\frac14-\tilde\lambda-3m^2-kp$. On each box $R$ is bounded by its second-order Taylor form
\[
R(x)\in R(x_c)+R'(x_c)h+R''(\mathrm{box})\,h^2/2 .
\]
Here $x_c=\tau_{\rm tip}-\tau_c$, and the profile derivatives come from the ODE ($k'=m$, $m'=p/4$, $p'=-kq$, $q'=1+kp$). The tip uncertainty enters only through $x_c$.

Together with the far-field stages of the earlier certificate, now run with the tight tip interval ($|\sigma_{\rm tip}-\mathrm{STIP}|\le5.5\cdot10^{-10}$), this gives
\[
\|(L-\tilde\lambda)\hat\psi_3\|_{L^2(\rho;[0,\infty))}\le1.2772\cdot10^{-6}\qquad(\text{previously }3.8051\cdot10^{-3})
\]
(\texttt{psi3res3.py}). The stage contributions are $5.7\cdot10^{-7}$ on $[\sigma_{\rm tip},8]$ and $1.14\cdot10^{-6}$ on $[8,32]$; the tail is negligible. The floor is now the fit of $\hat\psi_3$ itself and the chunk bounds of the profile on $[8,32]$.

In the triangular lemma, $\|g\|=\|(I-\Pi)\hat r\|/N'$ drops from $0.122$ to $4.1\cdot10^{-5}$. The rate $\omega'$ and the constant $C$ are unchanged, since $\eps_1\|\phi_3\|^2$ dominates them. In the nonlinear bootstrap the forcing $\alpha(I-\Pi)\hat r$ becomes negligible, so the $O(\alpha^2)$ forcing is now the main limit on the basin.

\section*{Reproducibility}
\texttt{run\_tight.sh} (bundle \texttt{s1\_tight.zip}) recompiles the CAPD programs and reruns every step in under a minute, given the step-1 folder (\texttt{tail.py}, \texttt{geometry.py}). The trust base is the same as for step 1.
\end{document}
